综述部分用于把案例研究放回更大的文献脉络中。

\subsection{职业社会学的贡献}

职业社会学强调职业路径并非个体单独选择的结果，而是制度、机会结构与社会资本共同作用的产物\cite{Super1957,Bourdieu1986}。这一视角帮助我们理解，为什么不同平台对同一艺人的价值并不恒定。

\subsection{媒体研究的贡献}

融合文化研究指出，新旧媒体并不是简单替代关系，而是在竞争与协作中重新分配注意力\cite{Jenkins2006}。因此，艺人的职业重启常常不是“离开旧平台”，而是通过新平台重新定义旧经验。

\subsection{形象管理研究的贡献}

形象管理研究强调公众人物需要在多重舞台之间持续完成自我呈现\cite{Goffman1959}。在危机情境中，这种呈现不再只是风格问题，而成为职业可持续性的重要组成部分。

\subsection{综述小结}

现有研究为本示例提供了三个关键启发：职业路径需要阶段化观察，平台迁移需要机制化解释，危机修复需要放到关系与叙事的双重层面讨论。
